%! Setting Up a Test for the Difference of Two Population Proportions — Unit 6, Lesson 6.10 — Warm-Up (Answer Key)
\documentclass[10pt]{article}
\usepackage{apstats-article}
\usepackage{apstats-key}   % pulls in -boxes; do NOT also load -boxes
\newcommand{\TallMath}[1]{$\displaystyle #1\rule[-1.4em]{0pt}{3.2em}$}

\begin{document}

\pageheader{Unit 6, Lesson 6.10}{Warm-Up --- Answer Key}
\namedateperiod

\vspace{0.15in}

\textbf{Spiral: Setting Up a One-Sample Test.}
A researcher claims that more than 40\% of adults in a city use a ride-share app weekly.
In a random sample of 250 adults, 115 reported doing so.

\begin{enumerate}
  \item \ansline{$H_0: p = 0.40$ (the proportion who use ride-share weekly is 40\%); $H_a: p > 0.40$ (more than 40\% use ride-share weekly). One-sided (right-tail) test.}

  \item Use \ans{$p_0$} (the null value, 0.40) for Large Counts in a one-sample test.

  \item $n \cdot p_0 = \ans{100}$, \quad $n(1-p_0) = \ans{150}$. \quad Both $\ge 10$? \ans{Yes}

  \item \ansline{No. Under $H_0: p_1=p_2$, both groups are assumed to share one unknown proportion. We pool the data to estimate that common proportion $\hat p_c$ and use it for all four Large Counts checks.}
\end{enumerate}

\end{document}
